
\label{sec:alpha-dos-vias-t1-rev8}

La clôture de \(\alpha\) ne commence ni par une constante métrologique ni par une
suite décimale recherchée. Elle commence dans le même état terminal enrichi qui
conserve la généalogie des trois canaux. Cette section présente dans une seule
chaîne causale deux voies vers l’échelle de \(\alpha\) et un prolongement long
de la retenue de la première. Le prolongement distant est noté
\(\mathsf T_{\alpha,\infty}\) ; il ne constitue ni une deuxième voie ni une troisième
dérivation.

\subsection{État terminal commun et 2 lectures internes coordonnées}

Soit
\[
 x_{\rm term}\in\mathcal X_{\rm TPK}^{\rm term,enr}
 \subseteq\mathcal X_{\rm TPK}^{[108]}.
\]
Le sélecteur terminal et la carte signée sont deux évaluations du même état :
\begin{equation}
 x_{\rm term}
 \xrightarrow{(\pi_{\rm term},R_{\rm sgn})}
 \bigl(B_{\rm term},U_{12}^{\rm sgn}\bigr),
 \label{eq:alpha-rev8-doble-lectura-terminal}
\end{equation}
avec
\[
 B_{\rm term}=
 \begin{pmatrix}
 021020\\001220\\100210
 \end{pmatrix}
\]
et
\[
 U_{12}^{\rm sgn}
 =(2378,1406,2479,-452,998,-551,-668,-204,-371,-322,-28,-997).
\]
La première lecture conserve les trois blocs terminaux sélectionnés ; la
seconde conserve la coordonnée entière signée nécessaire à la clôture
dodécaphasique. Il n’existe aucune factorisation de
\(B_{\rm term}\) vers \(U_{12}^{\rm sgn}\), et aucune n’est nécessaire. La dessiner effacerait précisément les
champs d’orientation, de retenue, de frontière et de registre qui distinguent les deux
lectures. À partir de la seconde, en revanche, les cartes réversibles sont bien construites
\[
 U_{12}^{\rm sgn}
 \xleftrightarrow{\ \mathcal H_{12}\ }K
 \xleftrightarrow{\ \mathcal E_{\rm dod}\ }
 (D_3K,D_4K,Q),
 \qquad
 \mathcal E_{\rm dod}=(D_3,D_4,Q).
\]
Le domaine de cette clôture est donc le profil enrichi du TPK ; une
trace brute de 108 positions est une projection par oubli et ne remplace pas ce
domaine.

\subsection{Voie A : clôture entière dodécaphasique}

Les évaluations archimédiennes des trois canaux et le sceau canonique sont
les vecteurs de douze blocs
\[
\begin{aligned}
P={}&(141,592,653,589,793,238,462,643,383,279,502,884),\\
E={}&(718,281,828,459,045,235,360,287,471,352,662,497),\\
\Phi={}&(618,033,988,749,894,848,204,586,834,365,638,117),\\
K={}&(234,543,140,729,659,824,621,058,914,794,146,601).
\end{aligned}
\]
Avec la frontière droite déclarée \(c_{13}=0\), la division euclidienne se
propage de droite à gauche :
\begin{equation}
\begin{aligned}
 s_m&=P_m+E_m-\Phi_m-K_m+c_{m+1},\\
 A_m&=s_m\bmod 1000,\\
 c_m&=\frac{s_m-A_m}{1000},
 \qquad m=12,11,\ldots,1.
\end{aligned}
\label{eq:alpha-rev8-via-entera}
\end{equation}
La sortie est
\[
 A=(007,297,352,569,283,800,997,285,105,472,380,663)
\]
et fixe le rationnel
\begin{equation}
 \alpha_{12}
 =0.007297352569283800997285105472380663.
 \label{eq:alpha-rev8-salida-corta}
\end{equation}
L’égalité abrégée \(P+E-\Phi-K=A\) doit se lire dans la carte de
concaténation entière. Si
\[
 \iota(v_1,\ldots,v_{12})
 =\sum_{m=1}^{12}v_m1000^{12-m},
\]
alors
\begin{equation}
 \boxed{\iota(P)+\iota(E)-\iota(\Phi)-\iota(K)=\iota(A)}.
 \label{eq:alpha-rev8-identidad-concatenada}
\end{equation}
En coordonnées, le cocycle ne peut être omis :
\[
 P+E-\Phi-K+C^+-1000C=A,
 \qquad C^+=(c_2,\ldots,c_{13}).
\]
La terminaison \(\ldots|380|663\) appartient à cette fenêtre finie et à sa
condition \(c_{13}=0\). Le théorème complet, le tableau des douze retenues et
l’inversion qui récupère \(K\) se trouvent dans le chapitre précédent ; ils ne sont pas
dupliqués ici.

Cette voie est un résultat interne exact avec une structure de départ explicite :
elle reçoit \(x_{\rm term}\), sa lecture signée, \(P,E,\Phi\), le sceau réversible
et la frontière. Elle ne reçoit ni \(\alpha\) ni CODATA en entrée. La comparaison
métrologique ne s’ouvre qu’après l’obtention de \eqref{eq:alpha-rev8-salida-corta}.

\subsection{Voie B : noyau de vacances \texorpdfstring{\(E_{21/22}\)}{E21/22}}

La seconde voie n’inverse pas les cartes de \(K\). Elle part de la densité exacte
de vacance
\[
 \varepsilon_{\rm v}=\log_{10}\!\left(\frac{10}{9}\right)
\]
et du noyau analytique déclaré
\begin{equation}
 \mathcal P(x)
 =2\pi x-\frac74x^2+\frac{x^3}{2\pi}
  +\frac{x^4}{20}-\frac{2x^5}{21}-\frac{x^6}{46}.
 \label{eq:alpha-rev8-kernel-e21-22}
\end{equation}
Le problème consiste à résoudre
\begin{equation}
 \mathcal P(x)=\log_{10}\!\left(\frac{10}{9}\right),
 \qquad 0<x<10^{-2}.
 \label{eq:alpha-rev8-ecuacion-e21-22}
\end{equation}
Le théorème de monotonie du noyau établit l’unicité de la racine positive
et l’encadrement
\[
 0.00729735256928379955
 <x_{21/22}<
 0.00729735256928379957.
\]
Son inverse est
\begin{equation}
 x_{21/22}^{-1}
 =137.0359990840000270200901282598749\ldots.
 \label{eq:alpha-rev8-inversa-e21-22}
\end{equation}
Par conséquent, le noyau complet \(E_{21/22}\), engendré par le second
lecteur du même état enrichi, détermine comme racine positive simple et
unique la même section limite que la voie dodécaphasique :
\[
 \boxed{\alpha_A=\alpha_B=\alpha_{\rm HMT}},
 \qquad
 \alpha_B=\operatorname*{arg\,zero}_{x\in I_\alpha}E_{21/22}(x).
\]
Le rationnel fini \(\alpha_{12}\) et la racine du jet \(\mathcal P=P_6\)
sont des projections finies des deux familles compatibles ; leur différence au-delà
de la résolution imprimée ne constitue pas une différence entre les deux
sections limites. Les intervalles de racine emboîtés et les carrés de
troncature commutatifs coordonnent les deux prolongements à toute profondeur.

Les six coefficients de \(P_6\) ne sont pas reconstruits à partir de la racine ni tirés
d’un tableau métrologique : ils sont le premier jet du noyau \(E_{21/22}\) produit
par le lecteur interne avec la généalogie \(120/45/11/495\) et le prolongement
\(T_{7:9}\). L’unicité de la racine et la dérivation interne des
coefficients sont des étapes distinctes d’une même voie B déjà close.

\subsection{Prolongement distant \texorpdfstring{\(\mathsf T_{\alpha,\infty}\)}{T alpha infini} de la voie entière}

\(\mathsf T_{\alpha,\infty}\) étend la récurrence de
\eqref{eq:alpha-rev8-via-entera} à une chaîne
de longueur \(N\). Pour éviter une ambiguïté d’indices, soit
\(\kappa(k)=1+((k-1)\bmod12)\). Avec une condition terminale à l’extrémité
distante, on calcule
\begin{equation}
\begin{aligned}
 s_k&=P_k+E_k-\Phi_k-K_{\kappa(k)}+c_{k+1},\\
 A_k&=s_k\bmod1000,\\
 c_k&=\left\lfloor\frac{s_k}{1000}\right\rfloor,
 \qquad k=N,N-1,\ldots,1.
\end{aligned}
\label{eq:alpha-rev8-t1}
\end{equation}
Le cas fini \(N=1000\) utilise mille triades, soit trois mille chiffres. La
retenue qui arrive de l’extrémité distante modifie le douzième
bloc visible :
\[
\begin{array}{rcl}
\text{fenêtre dodécaphasique courte, }c_{13}=0
 &:& \ldots|380|663,\\
\text{prolongement }\mathsf T_{\alpha,\infty}\text{ de mille triades}
 &:& \ldots|380|662|999\ldots.
\end{array}
\]
Ce ne sont pas deux valeurs rivales. Ce sont deux sémantiques de frontière de la même loi de
retenue : la première se ferme à la douzième position ; la seconde y transporte
des informations provenant d’une queue distante.

Le calcul reproductible vérifie exactement le mot long en utilisant
\(P,E,\Phi\) comme projections corrélées du même état coinductif
conjoint APP--TRIT--TPK qui publie également la coordonnée dodécaphasique
\(\alpha_{\mathrm{HMT}}\). Cette exécution certifie donc la propagation
de la retenue : elle ne place pas \(\pi,e,\varphi\) plus tôt dans la généalogie et n’ajoute pas
\(\alpha\) de l’extérieur. La projection conjointe
\[
 \mathcal G_9^{\mathrm{joint}}\longrightarrow
 \bigl(\pi,e,\varphi,\mathsf T_{\alpha,\infty}\bigr)
\]
se factorise dans le \cref{sec:alpha-t1-10000-rev9} : les trois canaux
archimédiens et le canal dodécaphasique s’ouvrent à partir du même état
enrichi ; on exécute ensuite \eqref{eq:alpha-rev8-t1} et l’on vérifie
l’indépendance du préfixe à l’égard des cinq conditions terminales de
retenue. La séparation des projections empêche la récurrence longue de
sélectionner rétrospectivement ses données initiales.

\subsection{Coordination causale de la clôture dodécaphasique et de la caractérisation \texorpdfstring{\(E_{21/22}\)}{E21/22} : domaines, codomaines et unicité}

La classification de ce bloc est fixée comme suit :
\begingroup
\raggedright
\begin{description}[style=nextline,leftmargin=3.2cm,labelwidth=2.8cm,labelsep=.4cm]
\item[Voie A : clôture dodécaphasique.]
Elle reçoit l’état terminal enrichi, la carte signée, le sceau, les trois
évaluations et la frontière explicite ; elle produit le rationnel
\eqref{eq:alpha-rev8-salida-corta} au moyen de l’identité entière avec retenue.

\item[Voie B : équation \(E_{21/22}\).]
Pour le noyau déclaré, la monotonie et le changement de signe déterminent une
unique racine dans \((0,10^{-2})\). La racine ne détermine pas rétrospectivement les
six coefficients du noyau.

\item[Prolongement \(\mathsf T_{\alpha,\infty}\).]
La récurrence longue reçoit les mille triades de référence, le sceau périodique
et la frontière distante, et transporte la retenue jusqu’à la fenêtre visible. Elle ne
constitue pas la voie B.

\item[Composition \(\mathcal G_9\to\mathsf T_{\alpha,\infty}\).]
L’arithmétique rationnelle conjointe produit dix mille chiffres, exécute la
récurrence de retenue et vérifie ensuite la stabilité du préfixe sous les
conditions terminales déclarées.
\end{description}
\endgroup

Le sélecteur terminal, les cartes réversibles et la récurrence courte sont
démontrés dans le \cref{chap:alpha} ; l’existence et l’unicité de la seconde voie
dans \eqref{eq:alpha-rev8-kernel-e21-22}--\eqref{eq:alpha-rev8-inversa-e21-22} ; et la composition longue dans le
\cref{sec:alpha-t1-10000-rev9}. Le supplément computationnel reproduit les
énumérations finies et les conditions terminales ; la démonstration repose
sur les identités et les récurrences exposées ici.

L’ordre causal public est ainsi clos :
\[
 x_{\rm term}
 \xrightarrow{(\pi_{\rm term},R_{\rm sgn})}
 (B_{\rm term},U_{12}^{\rm sgn})
 \longrightarrow
 \begin{cases}
 \text{voie A : clôture entière dodécaphasique},\\
 \text{voie B : racine du noyau }E_{21/22},
 \end{cases}
\]
avec \(\mathsf T_{\alpha,\infty}\) comme prolongement distant de la voie A et non
comme branche supplémentaire. Ce n’est
qu’après ces sorties que s’ouvre la reconnaissance métrologique externe.
